\section{Introduction}

When an autonomous drone completes a mission, the operator can say
what happened: the trajectory, the sensor readings, the commands
sent. What the operator \emph{cannot} say, without additional
machinery, is what \emph{would have happened} under different
conditions. Did the drone choose the safest route, or merely a safe
one? Would it have failed if the wind had been stronger? Was the
altitude within the envelope by design, or by luck?

These questions matter for certification. A regulator (EASA, FAA)
or an insurer does not only want to know that the drone was safe —
they want to know \emph{why} it was safe, and whether it would
remain safe under conditions that did not occur during the mission.
Static safety cases~\cite{kelly1998} cannot answer this: they describe
the intended behaviour, not the counterfactual.

Black-box logs (ArduPilot, PX4) record the actual mission, but they
do not replay it under modified conditions. Sensitivity analysis
techniques~\cite{saltelli2008} can explore the parameter space, but
they are offline and do not produce a verifiable link to the
recorded mission. Counterfactual reasoning~\cite{pearl2009} provides
the theoretical framework, but existing implementations do not
address the cryptographic requirements of an audit.

We present Counterfactual Safety Reasoning (CSR), a technique that
enables a drone to answer post-mission ``what if?' questions with
cryptographically verifiable answers. Given a recorded trajectory
and a hash chain of safety events (CertChain), CSR replays the
mission under modified conditions and produces a proof bound to the
original chain. A regulator can verify that ``if wind had been
$2\times$, the drone would have missed waypoint~3 by 40~m' without
trusting the operator.

CSR makes three contributions:

\begin{enumerate}
    \item A trajectory replay engine that re-runs the mission under
          modified conditions (wind scale, altitude offset, obstacle
          changes) and classifies the outcome (Better, Same, Worse,
          Failed).
    \item A proof construction that binds the counterfactual result
          to the original chain via SHA-256, so that tampering with
          either the chain or the result invalidates the proof.
    \item An evaluation on a 4-scenario case study (wind $0\times$,
          $1\times$, $20\times$, altitude $+150$~m), showing 4/4
          proofs verified and 3/3 attacks (tampered chain, tampered
          result, wrong key) rejected.
\end{enumerate}

We are explicit about limitations. The physics is synthetic
(1D/2D), the hash is FNV-1a (not SHA-256), and there are no formal
guarantees. CSR is a proof of concept, not a certified tool. But
it demonstrates a capability that no existing drone has: the ability
to prove, after the fact, what would have happened under different
conditions.

The paper is organised as follows. Section~II reviews safety cases
and cryptographic audit. Section~III presents the CSR design.
Section~IV describes the implementation. Section~V evaluates CSR
on the 4-scenario case study. Section~VI discusses the relation to
DO-326A and the ethical implications. Section~VII concludes.

% TODO (Andre):
% - Verify the claim "no existing drone has this" with a literature search
% - Add a footnote about FNV-1a vs SHA-256
% - Consider adding a sentence about the open-source availability
% - The ethical implications section needs more depth
